% Generated draft: requires mathematical and source review.
\documentclass[11pt]{article}
\usepackage[margin=25mm]{geometry}
\usepackage{fontspec}
\setmainfont{Latin Modern Roman}
\usepackage{amsmath,amssymb,mathtools,unicode-math}
\setmathfont{Latin Modern Math}
\usepackage{enumitem,needspace,xurl,hyperref}
\setlength{\parindent}{0pt}
\setlength{\parskip}{5pt}
\setlength{\emergencystretch}{3em}
\newcommand{\N}{\mathbb{N}}
\newcommand{\Z}{\mathbb{Z}}
\newcommand{\Q}{\mathbb{Q}}
\newcommand{\R}{\mathbb{R}}
\newcommand{\C}{\mathbb{C}}
\newcommand{\F}{\mathbb{F}}
\newcommand{\A}{\mathbb{A}}
\newcommand{\PP}{\mathbb{P}}
\newcommand{\E}{\mathbb{E}}
\newcommand{\eps}{\varepsilon}
\newcommand{\dd}{\,\mathrm{d}}
\begin{document}
\section*{UnsolvedMath ID: 800431}
\textbf{Source problem code:} catalogue entry
\subsection{Definitions and mathematical statement}
Let \(K\) be a number field, and let \(M_K\) denote its set of places. For \(v\in M_K\), write \(K_v\) for the completion of \(K\). Let \(\mathcal A_K\) be the adele ring of \(K\), namely the restricted product of the fields \(K_v\) with respect to their rings of integers at the nonarchimedean places. For a variety \(X/K\), write \(X(\mathcal A_K)\) for its adelic points, equipped with the usual adelic topology, and view \(X(K)\) diagonally inside this space. For \(S\subseteq X(K)\), interpret
\[
\mathcal C(S):=\overline{S}^{\,X(\mathcal A_K)}.
\]
Let \(f:X\to X\) be a morphism defined over \(K\), let \(Y\subset X\) be a subvariety defined over \(K\), and let \(P\in X(K)\). Define the forward orbit by
\[
\mathcal O_f(P):=\{f^n(P):n\in\mathbb Z_{\geq 0}\},\qquad f^0=\operatorname{id}_X.
\]
\textbf{Dynamical Brauer--Manin Conjecture, as supplied (Hsia--Silverman, 2009, [MR2537714]).} For every such \(K,X,f,Y,P\),
\[
\mathcal C\bigl(Y(K)\cap\mathcal O_f(P)\bigr)
=
Y(\mathcal A_K)\cap\mathcal C\bigl(\mathcal O_f(P)\bigr).
\]
Here the source expression \(Y\cap\mathcal O_f(P)\) is interpreted as \(Y(K)\cap\mathcal O_f(P)\), and both closures are taken in the same ambient space \(X(\mathcal A_K)\).

\textit{Unresolved qualifications.} The supplied record does not specify whether ``subvariety'' means a closed subvariety, nor does it detail its adelic topology convention. The displayed formulation uses the usual adelic topology and a common ambient closure. The closedness convention requires confirmation, since it affects the relationship between closure and membership in \(Y(\mathcal A_K)\). No properness hypothesis or additional restriction on the dynamics is supplied; none has been added. This is a statement-only transcription requiring human mathematical review.
\subsection{Short English statement}
The supplied conjecture says that the adelic closure of the orbit points lying on \(Y\) equals the intersection of \(Y\)'s adelic points with the adelic closure of the entire forward orbit. The registry labels this problem \emph{open}; that label is supplied metadata, not an independently verified assessment. The meaning of ``subvariety'' and the adelic closure convention require review.
\subsection{Sources}
\begin{itemize}
\item Local UnsolvedMath registry, permanent ID \(800431\), exact source problem code \texttt{catalogue entry}; supplied complete source record, upstream revision \texttt{d371ee917ac9647fbb06cd6fc529deec6f3dda82}. The record supplies the status \emph{open} and was updated on 2026-10-16.
\item Benedetto, DeMarco, Ingram, Jones et al., \emph{Current Trends and Open Problems in Arithmetic Dynamics} (2018), Conjecture 54, identified in the supplied record as the Dynamical Brauer--Manin Conjecture. Supplied source URL: \url{https://arxiv.org/abs/1806.04980}. The record reports access on 2026-10-16 and extraction by \texttt{arxiv\_latex}; these are supplied provenance notes, not a claim of independent source access or verification.
\item The supplied statement attributes the conjecture to Hsia--Silverman (2009), with reference label [MR2537714]. No full bibliographic entry for that reference is supplied. The registry's proposer field lists Benedetto, DeMarco, Ingram, Jones et al.; the statement's Hsia--Silverman attribution is preserved separately.
\end{itemize}
\end{document}
